\chapter{EXPERIMENT REQUIREMENTS AND CONSTRAINTS}
\label{sec:requirements}
\section{Functional Requirements}

\begin{longtable}{|p{2cm}|p{12cm}|}
\caption{Functional Requirements} \\
\hline 
\textbf{ID} & \textbf{Requirement text} \\
\hline 
\endfirsthead
\caption[]{Functional Requirements -- \textit{Continued from previous page}} \\
\hline
\textbf{ID} & \textbf{Requirement text} \\
\hline
\endhead
\hline
\multicolumn{2}{r}{\textit{Continued on next page}} \\
\endfoot
\hline
\endlastfoot
F.1. & The experiment shall measure the atmospheric CH\textsubscript{4}, CO\textsubscript{2} and H\textsubscript{2}O concentrations of ambient air simultaneously with an off-the-shelf \gls{ndir} sensor. \\
\hline
F.2. & The measurement chamber shall be pressurised with ambient air. \\
\hline
F.3. & The air in the measurement chamber shall be exchanged. \\
\hline
F.4. & The experiment shall measure the pressure inside of the measurement chamber. \\
\hline
F.5. & The experiment shall measure the temperature inside of the measurement chamber. \\
\hline
F.6. & The experiment shall stabilize the temperature inside of the measurement chamber. \\
\hline
F.7. & The experiment shall measure the humidity inside of the measurement chamber. \\
\hline
F.8. & The experiment shall prevent the condensation of water inside of the measurement chamber. \\
\hline
F.9. & The experiment shall protect electronics from undercooling or overheating. \\
\hline
F.10. & The experiment should measure ambient temperature outside. \\
\hline
F.11. & The experiment should measure ambient pressure outside. \\
\hline
F.12. & The experiment should measure ambient humidity outside. \\
\hline
F.13. & The experiment shall send data via E-Link. \\
\hline
F.14. & The experiment shall receive data via E-Link. \\
\end{longtable}

%%%%%%%%%%%%%%%%%%%%%%%%%%%%%%%%%%%%%%%%%%%%%%%%%%%%%%%%%%%%%%%%%%%%%%%%%%%%%%%
\section{Performance Requirements}

\begin{longtable}{|p{2cm}|p{12cm}|}
\caption{Performance Requirements}
\label{tab:performance_req}\\
\hline
\textbf{ID} & \textbf{Requirement text} \\
\hline
\endfirsthead
\caption[]{Performance Requirements -- \textit{Continued from previous page}} \\
\hline
\textbf{ID} & \textbf{Requirement text} \\
\hline
\endhead
\hline
\multicolumn{2}{r}{\textit{Continued on next page}} \\
\endfoot
\hline
\endlastfoot
P.1.1 & The \gls{ndir} sensor shall measure the concentration of methane with a range of 0.4 to 3.0 ppm. \\
\hline
P.1.2 & The \gls{ndir} sensor shall measure the concentration of methane with an uncertainty of 0.5 ppm or less. \\
\hline
P.2.1 & The \gls{ndir} sensor shall measure the concentration of carbondioxide with a range of 200 to 500 ppm. \\
\hline
P.2.2 & The \gls{ndir} sensor shall measure the concentration of carbondioxide with an uncertainty of 5 ppm or less. \\
\hline
P.3.1 & The \gls{ndir} sensor shall measure the concentration of water with a range of 1 to 30,000 ppm. \\
\hline
P.3.2 & The \gls{ndir} sensor shall measure the concentration of water with an uncertainty of 5 ppm or less. \\
\hline
%P.4.1 & The \gls{ndir} sensor shall have a measurement time for the CH\textsubscript{4}, CO\textsubscript{2} and H\textsubscript{2}O concentrations of 20 s during ascent and descent. \\
%\hline
%P.4.2 & The \gls{ndir} sensor shall have a measurement time for the CH\textsubscript{4}, CO\textsubscript{2} and H\textsubscript{2}O concentrations of 50 s during the floating phase. \\
%\hline
P.4.1 & The pressure inside of the measurement chamber shall stay within a range of 2.5 and 3.2 bar. \\
\hline
P.4.2 & The pressure inside of the measurement chamber should stay within a range of 2.8 and 3 bar. \\
\hline
P.5 & The pressurisation system shall have an air throughput of at least 1 l/min.\\
\hline
P.6.1 & The experiment shall measure the pressure inside of the measurement chamber in a range of 0.3 to 4 bar. \\
\hline
P.6.2 & The experiment shall measure the pressure inside of the measurement chamber with a maximum uncertainty of 0.02 bar. \\
\hline
P.7 & The temperature of the air inside of the measurement chamber shall stay within a range of 15 and 40\degree C while active. \\
\hline
P.8 & The temperature of the air inside of the measurement chamber should stay within a range of 19 and 24\degree C while active. \\
\hline
P.9.1 & The experiment shall measure the temperature inside of the measurement chamber in a range of -40 to +85 \degree C. \\
\hline
P.9.2 & The experiment shall measure the temperature inside of the measurement chamber with a maximum uncertainty of 1.5\degree C. \\
\hline
P.10 & The experiment shall keep the humidity inside of the measurement chamber below 90\% r.H. \\
\hline
P.11.1 & The experiment shall measure the humidity inside of the measurement chamber in a range of 0 to 100\% r.H. \\
\hline
P.11.2 & The experiment shall measure the humidity inside of the measurement chamber with a maximum uncertainty of 3\% r.H. \\
\hline
P.12 & The experiment should measure the temperature, pressure and humidity inside of the measurement chamber at a rate over 10 Hz. \\
\hline
P.13.1 & The temperature in the electronic compartment shall stay within a range
of -25 and 85°C. \\
\hline
P.13.2 & The temperature in the electronic compartment should stay within a range of -15 and 60°C. \\
\hline
P.14.1 & The experiment should measure the ambient pressure outside in a range of 20 to 1100 mbar. \\
\hline
P.14.2 & The experiment should measure the ambient pressure outside with a maximum uncertainty of 20 mbar. \\
\hline
P.15.1 & The experiment should measure the ambient temperature outside in a range of -70 to +60 \degree C. \\
\hline
P.15.2 & The experiment should measure the ambient temperature outside with a maximum uncertainty of 1.5\degree C. \\
\hline
P.16.1 & The experiment should measure the ambient humidity outside in a range of 0 to 100\% r.H. \\
\hline
P.16.2 & The experiment should measure the ambient humidity outside with a maximum uncertainty of 3\% r.H. \\
\hline
P.17 & The experiment should measure the ambient temperature, pressure and humidity outside at a rate over 10 Hz. \\
\hline
P.18.1 & The experiment shall operate in the shock profile of the BEXUS vehicle flight and launch. \\
\hline
P.18.2 & The \gls{ndir} sensor should withstand the shocks of the landing. \\
\hline
P.19 & The data storage shall be able to store at least 2GB. \\
\hline
P.20 & The experiment should receive data via E-Link with an uplink rate under 10 kbps. \\
\hline
P.21 & The experiment shall send the obtained data with the E-Link with a downlink rate under 100 kbps. \\
\end{longtable}

%%%%%%%%%%%%%%%%%%%%%%%%%%%%%%%%%%%%%%%%%%%%%%%%%%%%%%%%%%%%%%%%%%%%%%%%%%%%%%%
\section{Design Requirements}

\begin{longtable}{|p{2cm}|p{12cm}|}
\caption{Design Requirements} \\
\hline
\textbf{ID} & \textbf{Requirement text} \\
\hline
\endfirsthead
\caption[]{Design Requirements -- \textit{Continued from previous page}} \\
\label{tab:design_req}\\
\hline
\textbf{ID} & \textbf{Requirement text} \\
\hline
\endhead
\hline
\multicolumn{2}{r}{\textit{Continued on next page}} \\
\endfoot
\hline
\endlastfoot
D.1 & The experiment shall operate in the temperature profile of the BEXUS vehicle flight and launch. \\
\hline
D.2.1 & The experiment shall operate in the pressure profile of the BEXUS vehicle flight and launch. \\
\hline
D.2.2 & The measurement chamber shall withstand a pressure difference of up to 4 bar. \\
\hline
D.3 & The experiment shall operate in the vibration profile of the BEXUS vehicle flight and launch. \\
\hline
D.4 & The experiment power connector shall be an Amphenol PT02E8-4P, Code A.\\
\hline
D.5 & The experiment data connector shall be an Amphenol RJF21B connector. \\
\hline
D.6 & The experiment shall not disturb or harm the launch vehicle. \\
\hline
D.7 & The experiment shall consume less than 150 Wh during the flight. \\
\hline
D.8 & The experiment shall operate under 28.8\,V. \\
\hline
D.9.1 & The experiment shall have no currents exceeding 2.4 A. \\
\hline
D.9.2 & The experiment should have no currents exceeding 1.4 A. \\
\hline
D.10 & The experiment should stay within a size limit of 15x20x20 cm\textsuperscript{3}. \\
\hline
D.11 & The experiment shall stay within a weight limit of 5 kg. \\
\hline
D.12 & The materials of the experiment shall not outgas. \\
\hline
D.13 & The experiment shall have access to fresh air from outside. \\
\hline
D.14 & The data storage shall be protected from water due to condensation or during landing. \\
\hline
D.15 & The budget for all components except the K96 sensor should stay below 2000\texteuro. \\
\end{longtable}

%%%%%%%%%%%%%%%%%%%%%%%%%%%%%%%%%%%%%%%%%%%%%%%%%%%%%%%%%%%%%%%%%%%%%%%%%%%%%%%
\section{Operational Requirements}

\begin{longtable}{|p{2cm}|p{12cm}|}
\caption{Operational Requirements} \\
\hline
\textbf{ID} & \textbf{Requirement text} \\
\hline
\endfirsthead
\caption[]{Operational Requirements -- \textit{Continued from previous page}} \\
\hline
\textbf{ID} & \textbf{Requirement text} \\
\hline
\endhead
\hline
\multicolumn{2}{r}{\textit{Continued on next page}} \\
\endfoot
\hline
\endlastfoot
O.1 & The experiment shall be able to conduct measurements autonomously in case connection with the ground segment is lost. \\
\hline
O.2 & The experiment shall accept a request for radio silence at any time while on the launch pad. \\
\hline
O.3 & The experiment shall store all data on board and continuously send data to ground via E-Link. \\
\hline
O.4 & The experiment shall automatically register and respond to high r.H. inside the chamber in order to lower the humidity. \\
\end{longtable}

%%%%%%%%%%%%%%%%%%%%%%%%%%%%%%%%%%%%%%%%%%%%%%%%%%%%%%%%%%%%%%%%%%%%%%%%%%%%%%%
\section{Constraints}

\begin{longtable}{|p{2cm}|p{12cm}|}
\caption{Constraints} \\
\hline
\textbf{ID} & \textbf{Requirement text} \\
\hline
\endfirsthead
\caption[]{Constraints -- \textit{Continued from previous page}} \\
\hline
\textbf{ID} & \textbf{Requirement text} \\
\hline
\endhead
\hline
\multicolumn{2}{r}{\textit{Continued on next page}} \\
\endfoot
\hline
\endlastfoot
C.1 & The thermal chamber at the \gls{irf} goes only down to -40\degree C. \\
\end{longtable}